\documentclass[11pt]{article}
\usepackage[a4paper,margin=1in]{geometry}
\usepackage[T1]{fontenc}
\usepackage{lmodern,microtype}
\usepackage{amsmath,amssymb,amsthm,mathtools}
\usepackage{booktabs,enumitem,xcolor}
\usepackage[numbers,sort&compress]{natbib}
\usepackage[colorlinks=true,linkcolor=blue!55!black,citecolor=blue!55!black,urlcolor=blue!55!black]{hyperref}
\linespread{1.06}
\newtheorem{theorem}{Theorem}[section]
\newtheorem{proposition}[theorem]{Proposition}
\newtheorem{corollary}[theorem]{Corollary}
\theoremstyle{remark}\newtheorem{remark}[theorem]{Remark}
\newcommand{\C}{\mathbb C}
\newcommand{\Prob}{\mathcal P}

\title{Global Cloud Gauge and Empirical-Root Tightness\\
\large Exact Moment Compactness for Rank-Growing Conjugate Spectra}
\author{Liang Wang}
\date{July 2026}

\begin{document}
\maketitle

\begin{abstract}
A rank-growing spectral cloud requires one gauge for the whole cloud, not a
separate normalization for each quartet or shell. For any nonconstant finite
multiset $\Lambda$, we remove its barycenter $m_\Lambda$ and divide by its
centered RMS radius $s_\Lambda$. The associated empirical probability
measure then has exact complex mean zero and exact second absolute moment one.

This elementary identity has a strong uniform consequence. For every family
of globally normalized clouds,
\[
 \mu_\Lambda(\{|z|>R\})\le R^{-2}.
\]
Hence the empirical measures are uniformly tight on $\C$ and every sequence
has a weakly convergent subsequence. The result is exact and independent of
the ranks, shell geometry, or noise schedule.

For the 32 RH-222 clouds, direct recomputation agrees with the archived gauge
to machine precision. The maximum mean residual is
$4.99\times10^{-17}$, the maximum second-moment error is
$6.67\times10^{-16}$, the maximum fourth moment is $2.31691$, and the largest
normalized modulus is $1.83188$. All tested tail inequalities have positive
slack.

Tightness provides macroscopic probability compactness. It neither selects a
unique weak limit nor proves local finiteness of the unweighted root divisor.
That distinction becomes the next exact obstruction.
\end{abstract}

\section{One gauge per growing cloud}

RH-220 completed the center-radius-shape dictionary for one quartet, and
RH-222 extends the root set to ranks $4$ through $35$
\citep{WangRH222}. Applying a separate gauge to each shell would erase their
relative positions. We instead normalize the entire selected multiset at
once.

Let
\[
 \Lambda=\{\lambda_1,\ldots,\lambda_n\}\subset\C,\qquad n\ge2,
\]
with not all roots equal. Define
\begin{equation}\label{eq:gauge}
 m_\Lambda=\frac1n\sum_{j=1}^n\lambda_j,\qquad
 s_\Lambda=\left(\frac1n\sum_{j=1}^n
 |\lambda_j-m_\Lambda|^2\right)^{1/2}>0,
\end{equation}
and normalized roots
\begin{equation}
 q_j=\frac{\lambda_j-m_\Lambda}{s_\Lambda}.
\end{equation}
The empirical measure is
\begin{equation}\label{eq:empirical}
 \mu_\Lambda=\frac1n\sum_{j=1}^n\delta_{q_j}.
\end{equation}

\section{Exact moment identities}

\begin{proposition}[Global gauge identities]\label{prop:moments}
Every globally normalized cloud satisfies
\begin{equation}\label{eq:identities}
 \int_\C z\,d\mu_\Lambda(z)=0,\qquad
 \int_\C |z|^2\,d\mu_\Lambda(z)=1.
\end{equation}
If $\Lambda$ is conjugate closed, then $\mu_\Lambda$ is invariant under
complex conjugation and $m_\Lambda\in\mathbb R$.
\end{proposition}
\begin{proof}
The first identity is
\[
 \frac1{ns_\Lambda}\sum_j(\lambda_j-m_\Lambda)=0.
\]
The second is exactly the definition of $s_\Lambda$. Conjugation permutes
the roots and therefore the atoms; it also fixes their mean.
\end{proof}

These identities do not depend on the root ordering. They remain valid when
real singleton shells make the rank odd.

\section{Uniform tightness}

\begin{theorem}[Second-moment tightness]\label{thm:tight}
Let $\{\Lambda_\alpha\}$ be any family of nonconstant finite clouds, each
normalized by \eqref{eq:gauge}. Then
\begin{equation}\label{eq:tail}
 \sup_\alpha\mu_{\Lambda_\alpha}(\{|z|>R\})
 \le \min(1,R^{-2}),\qquad R>0.
\end{equation}
Consequently the family
$\{\mu_{\Lambda_\alpha}\}\subset\Prob(\C)$ is uniformly tight.
\end{theorem}
\begin{proof}
On $\{|z|>R\}$, one has $1\le |z|^2/R^2$. Therefore
\[
 \mu_{\Lambda_\alpha}(|z|>R)
 \le R^{-2}\int |z|^2\,d\mu_{\Lambda_\alpha}=R^{-2}
\]
by Proposition~\ref{prop:moments}. For every $\varepsilon>0$, choosing
$R\ge\varepsilon^{-1/2}$ leaves mass at most $\varepsilon$ outside the compact
closed disk.
\end{proof}

\begin{corollary}[Weak subsequential compactness]\label{cor:prok}
Every sequence of globally normalized empirical cloud measures has a weakly
convergent subsequence in $\Prob(\C)$.
\end{corollary}
\begin{proof}
The complex plane is a Polish space. Apply Prokhorov's theorem to
Theorem~\ref{thm:tight}; see \citet{Billingsley1999}.
\end{proof}

This closes the tightness item in the RH-221 roadmap at the level of empirical
probability measures. No numerical asymptotic fit is needed.

\section{What weak convergence preserves}

If $\mu_{n_j}\Rightarrow\mu$, conjugation symmetry passes to the limit because
bounded continuous test functions can be conjugated. The mean and second
moment require care: $z$ and $|z|^2$ are unbounded. Portmanteau gives only
\begin{equation}
 \int |z|^2\,d\mu(z)\le1.
\end{equation}
Equality needs uniform integrability of $|z|^2$. A uniform $(2+\epsilon)$
moment bound would suffice, but the finite maximum fourth moment reported
below is not an all-level theorem.

Thus one must not infer that every subsequential limit automatically retains
the complete RMS gauge. Mass can vanish to infinity while carrying a
nonzero fraction of the second moment, even though probability tightness
holds.

\section{Physical audit}

For each of the 32 endpoints, the raw selected roots are normalized again,
without reading the archived normalized array. The two representations are
then compared. Tail masses are evaluated at
\[
 R=1,\;1.25,\;1.5,\;2,\;3
\]
against the exact bound \eqref{eq:tail}.

\begin{center}
\begin{tabular}{lr}
\toprule
diagnostic & maximum or minimum\\
\midrule
absolute empirical mean & $4.99\times10^{-17}$\\
second-moment error & $6.67\times10^{-16}$\\
fourth absolute moment & $2.31691$\\
normalized root modulus & $1.83188$\\
minimum tail-bound slack & $0.11111$\\
\bottomrule
\end{tabular}
\end{center}

In particular every audited normalized root lies inside $|z|<2$. This is
stronger than the universal $75\%$ mass guarantee at $R=2$, but it is only a
finite observation.

For reference, the theorem supplies scale-independent certificate radii
\[
 R_{0.25}=2,\quad R_{0.10}=\sqrt{10},\quad
 R_{0.05}=\sqrt{20},\quad R_{0.01}=10.
\]
These values are deliberately not tuned to the observed maximum.

\section{Probability mass versus divisor mass}

The empirical measure gives each root weight $1/n$. A holomorphic zero
divisor gives each simple root integer weight one. Therefore tightness of
$\mu_\Lambda$ says that a fixed compact set can contain a fixed
\emph{fraction} of an expanding cloud. Multiplying that fraction by
$n\to\infty$ may force the integer divisor mass on the same compact set to
diverge.

This is not a technicality. Probability-measure compactness and local
finiteness of zero divisors are different topologies. The next paper proves
their direct incompatibility for a tight rank-growing cloud. The likely
escape is to use reciprocal resonances, as in the fixed-noise regularized
Fredholm determinant of RH-7 \citep{WangRH7}, rather than the resonances
themselves.

\section{Claim boundary}

The exact achievements are:
\begin{enumerate}[leftmargin=2em]
\item one affine gauge for each full cloud;
\item zero mean and unit second moment;
\item uniform tightness and weak subsequential compactness.
\end{enumerate}
No unique weak measure, all-level fourth-moment bound, unweighted divisor
limit, or locally uniform determinant is proved. Gate A remains open, and no
later Hilbert--Polya or arithmetic claim is made.

\bibliographystyle{plainnat}
\bibliography{references}
\end{document}
